\section{从五级流水线走向超标量前端}

五级流水线默认每拍取一条、译一条、发一条。现代核心要把每拍退休指令数继续提高，首先要扩大前端宽度：取指单元按字节块读取指令，分支预测器给出下一取指地址，边界识别逻辑从字节流中切出若干条指令，译码器再把它们变成内部操作。RISC 定长指令的边界预先确定；x86 变长指令必须先识别前缀、opcode、ModR/M、SIB、位移和立即数，复杂指令还可能拆成多个微操作。

“四发射”不是保证每拍总能执行四条，而是规定前端、重命名和提交等位置最多能处理四个操作。若取指块在 Cache line 末尾断开、预测目标未命中、译码遇到复杂指令，或者后端队列没有空位，当拍有效宽度都会下降。设计者因此要给每个边界同时规定宽度和背压条件。

\begin{longtable}{L{0.18\linewidth}L{0.25\linewidth}L{0.25\linewidth}L{0.22\linewidth}}
\toprule
前端边界 & 保存的状态 & 正常推进条件 & 停止或清空条件 \\
\midrule
取指队列 & PC、字节块、预测信息 & 指令 Cache 返回且队列有空位 & I-Cache miss、队列满或重定向 \\
译码队列 & 指令边界、操作数编码、异常标记 & 后端可接受对应数量的微操作 & 复杂译码占用、重命名端背压 \\
重命名端 & 架构源/目的、物理映射、ROB 序号 & 物理寄存器、ROB 和发射队列都有空位 & 任一资源不足时整组停顿 \\
恢复检查点 & 预测分支前的映射表与队列位置 & 分支按预测方向执行 & 误预测或更早异常触发恢复 \\
\bottomrule
\end{longtable}

同一拍进入的多条指令还存在组内依赖。若第 1 条写 \code{x1}，第 3 条读 \code{x1}，第 3 条不能读取本拍开始前的旧映射，而要读取第 1 条刚分配的新物理寄存器标签。重命名逻辑因此不是四份互不相干的查表器，而是一条按程序顺序组合展开的映射链；宽度越大，这条链的比较、旁路和空闲寄存器分配越重。

\topic{寄存器重命名把名字依赖变成数据标签}

架构寄存器数量由 ISA 固定，但核心内部可以准备更多物理寄存器。重命名表 RAT（register alias table）记录“当前最新的架构寄存器值在哪个物理寄存器”；空闲表记录尚未分配的物理寄存器；ROB（reorder buffer，重排序缓冲）按程序顺序保存每条在途操作的目的映射、完成状态和异常状态。

考虑下面四条指令，初始映射为 \code{x1->p1}、\code{x2->p2}、\code{x3->p3}：

\begin{longtable}{L{0.09\linewidth}L{0.28\linewidth}L{0.20\linewidth}L{0.20\linewidth}L{0.13\linewidth}}
\toprule
指令 & 架构动作 & 重命名后的源 & 新目的 & 被替代映射 \\
\midrule
I1 & \code{x1=x2+x3} & \code{p2,p3} & \code{x1->p32} & \code{p1} \\
I2 & \code{x4=x1-1} & \code{p32} & \code{x4->p33} & x4 的旧映射 \\
I3 & \code{x6=x7*x8} & x7、x8 的当前映射 & \code{x6->p34} & x6 的旧映射 \\
I4 & \code{x1=x5+2} & x5 的当前映射 & \code{x1->p35} & \code{p32} \\
\bottomrule
\end{longtable}

I2 对 I1 的 RAW（read after write，先写后读）是真数据依赖，重命名后仍由标签 \code{p32} 保留。I4 再写 \code{x1} 形成 WAW 名字依赖，但它得到 \code{p35}，不会覆盖 I1 尚未被 I2 使用的 \code{p32}。类似地，一条较晚指令写某个寄存器，不会破坏较早指令对旧物理值的读取，因此 WAR 名字依赖也被消除。

旧物理寄存器不能在重命名时立即回收。I4 建立 \code{x1->p35} 后，\code{p32} 仍可能被 I2 或其他更早的在途指令读取；只有 I4 按序退休，确认所有更早指令都不会回滚时，才可把它替代的映射放回空闲表。这个回收边界是重命名正确性的关键：分配发生在推测路径上，释放必须服从不可回退的提交顺序。

\topic{唤醒、选择与执行端口}

重命名后的操作进入发射队列。每个源操作数要么已经有值，要么保存生产者的物理标签。执行单元完成结果后发布“标签+值”或至少发布完成标签；发射队列比较所有等待源，命中的项把对应源标成就绪。这一步称为 wakeup。随后 select 逻辑从全部源已就绪的操作中选择可用执行端口，这一步还要满足端口类型、年龄、公平性和每拍发射宽度。

假设 I1 是 1 拍加法，I3 是 3 拍乘法，I2 等待 I1，I4 与它们无关。一个可能的时间线如下：

\begin{longtable}{L{0.12\linewidth}L{0.20\linewidth}L{0.20\linewidth}L{0.20\linewidth}L{0.18\linewidth}}
\toprule
拍 & ALU0 & ALU1 & 乘法器 & 队列与提交状态 \\
\midrule
1 & I1 发射 & I4 发射 & I3 发射 & I2 等待 \code{p32} \\
2 & I2 被唤醒并发射 & 空闲 & I3 第 2 拍 & I1、I4 已执行完成 \\
3 & 空闲 & 空闲 & I3 第 3 拍 & I2 已执行完成 \\
4 & 空闲 & 空闲 & I3 返回结果 & 全部执行完成，仍按 ROB 顺序退休 \\
\bottomrule
\end{longtable}

“较晚指令先执行完成”不等于“较晚指令先提交”。I4 在第 2 拍已经算完，但若它前面的 I3 尚未完成，I4 的结果只能留在物理寄存器与 ROB 状态中。发射队列按就绪关系选择，ROB 按程序顺序退休；这两张表分别回答“现在能算谁”和“现在能让谁对软件可见”。

端口冲突属于结构冒险。两条源操作数都就绪的乘法若只有一个乘法端口，只能选择一条；未选中的操作保持在队列中，不能丢失就绪状态。若执行单元返回端口数量少于同拍完成数量，还要提前调度或在返回侧排队，否则“都算完了”仍可能因结果总线冲突而不能写入物理寄存器。

\topic{Load/Store Queue 维护地址、数据与顺序}

普通整数操作只依赖寄存器标签，内存操作还依赖地址和内存顺序。LSQ（load/store queue）通常按程序顺序保存加载项与存储项。加载项记录有效地址、访问宽度、目标物理寄存器和完成状态；存储项记录地址、数据、字节掩码和是否已经退休。地址生成单元可以乱序算出地址，但存储通常要等到退休后才进入对 Cache 可见的 store buffer，避免错误路径留下不可撤销写入。

一条 load 在执行前至少要检查三类更早 store：

\begin{longtable}{L{0.24\linewidth}L{0.32\linewidth}L{0.34\linewidth}}
\toprule
更早 store 的状态 & load 的处理 & 原因 \\
\midrule
地址已知且不重叠 & 可越过该 store 访问 Cache & 已证明不存在内存依赖 \\
地址和数据已知且重叠 & 从 store queue 转发匹配字节 & 最新值尚未进入 Cache \\
地址未知 & 等待，或推测越过并保留重放条件 & 可能存在尚未发现的同址写 \\
地址重叠但数据未就绪 & load 等待数据 & 不能从 Cache 读取旧值 \\
\bottomrule
\end{longtable}

推测越过未知地址能提高性能，但必须可验证。设较早的 S1 地址计算较慢，较晚的 L2 先从地址 \code{0x1000} 读取旧值；随后 S1 算出自己的地址也是 \code{0x1000}。LSQ 比较发现 L2 违反了“应观察最近更早 store”的规则，必须使 L2 及其依赖操作失效，从 L2 或更早的安全边界重放。若核心只有推测而没有冲突检测与重放，偶发的数据相关就会变成不可重复的软件错误。

跨页、未对齐和部分重叠使比较更复杂。一次 8-byte load 可能跨两个 Cache line 或两个页面；一个 4-byte store 也可能只覆盖其中部分字节。LSQ 必须按物理地址和字节掩码组合多个来源，并保证异常检查先于不可撤销的部分结果。MMIO 访问通常不允许使用普通 load 的推测、合并和重放规则，需要等更早操作达到规定边界后再按设备属性发出。

\section{ROB 把乱序完成重新排列为按序退休}

ROB 项按程序顺序分配，常见字段包括操作类型、目的架构寄存器、当前物理寄存器、被替代的旧物理寄存器、完成位、异常原因和分支预测信息。执行单元只更新对应 ROB 项的完成或异常状态；提交端只检查队首若干项。队首操作完成且无异常时，提交端才更新架构恢复所需状态、允许 store 进入 store buffer，并回收旧物理寄存器。

\begin{longtable}{L{0.19\linewidth}L{0.27\linewidth}L{0.25\linewidth}L{0.19\linewidth}}
\toprule
边界 & 可以发生的动作 & 不能据此推断的动作 & 失败处理 \\
\midrule
执行完成 & 物理结果写回，唤醒依赖者 & 结果已经成为不可回退状态 & 在 ROB 中记录异常或重放 \\
ROB 队首完成 & 按退休宽度检查提交资格 & 较晚项也可越过异常提交 & 较早异常阻止本项退休 \\
指令退休 & 架构映射确定，旧物理寄存器可回收 & store 已经到达内存介质 & store 只获得进入写缓冲资格 \\
store buffer 接受 & 写获得一致性与 Cache 路径资源 & 其他核心已经观察到写 & 由内存序和一致性继续约束 \\
\bottomrule
\end{longtable}

精确异常依靠 ROB 队首实现。若 I3 除零，而 I4、I5 已在错误路径或推测路径上执行完成，提交端先让 I1、I2 正常退休；看到队首 I3 的异常后停止提交，恢复到 I3 之前的架构映射，清除 I3 及其后所有在途操作，再跳转到异常入口。软件看到的状态与顺序机器完全一致：I1、I2 已发生，I3 及以后均未发生。

\topic{分支误预测恢复的是一组状态，不只是 PC}

预测分支进入前端时，后续指令已使用预测 PC 取指并修改推测状态。分支执行后若方向或目标错误，核心至少要恢复取指 PC、RAT、空闲表位置、ROB 尾指针、发射队列和 LSQ 中属于错误路径的项。只改 PC 会留下错误路径分配的物理寄存器和队列项，最终造成资源泄漏或错误依赖。

一种做法是在每个未决分支处分配检查点，保存 RAT 快照和若干队列尾指针。误预测时恢复最近对应检查点，丢弃年轻项，前端从正确目标重新取指。另一种做法利用 ROB 中按顺序保存的旧映射逐项回滚，硬件较省但恢复拍数随错误路径长度增加。检查点数量也有限；未决分支过多时，前端即使预测仍有空间，也可能因没有恢复资源而停顿。

分支恢复与异常恢复的共同原则是“最早失败者优先”。同一窗口里可能同时报告一个较晚分支误预测和一个更早 load 的页错误；核心必须按程序顺序选择最早需要重定向的事件。否则先按较晚分支恢复，随后又发现更早异常，会让映射表和队列状态经历无法证明正确的组合。

\section*{章末练习}

\begin{chapterexercise}{乱序重命名}
初始 \code{x1->p1}、\code{x2->p2}。依次重命名 \code{x1=x2+1}、\code{x3=x1+2}、\code{x1=x4+3}，可用物理寄存器从 \code{p20} 开始。写出三条指令的源标签、新目的标签和每次被替代的旧映射。
\exercisecheck{学习目标}{区分 RAW 真依赖与 WAW/WAR 名字依赖}
\end{chapterexercise}

\begin{chapterexercise}{唤醒与端口}
两条乘法的源都在同拍就绪，但核心只有一个乘法端口。说明 select 逻辑如何处理未选中的一条，以及为什么不能清除它的就绪位。
\exercisecheck{学习目标}{就绪不等于已经获得执行资源}
\end{chapterexercise}

\begin{chapterexercise}{存储转发}
较早的 8-byte store 写 \code{0x1000--0x1007}，较晚的 4-byte load 读 \code{0x1004--0x1007}。说明 load 应从哪里取值；若 store 地址在 load 完成后才算出，应怎样恢复。
\exercisecheck{学习目标}{按物理地址与字节掩码检查依赖，并能重放违规 load}
\end{chapterexercise}

\begin{chapterexercise}{精确异常}
I1、I2 已完成，I3 报缺页，I4、I5 已执行完成但尚未退休。写出 ROB 提交端从发现 I3 到进入异常处理程序的状态变化。
\exercisecheck{学习目标}{执行完成与架构提交必须分离}
\end{chapterexercise}

\begin{chapterexercise}{误预测恢复}
列出分支误预测时除 PC 外至少四类必须恢复或清除的状态，并解释遗漏物理寄存器空闲表会产生什么后果。
\exercisecheck{学习目标}{恢复对象包含映射、队列、所有权和资源计数}
\end{chapterexercise}

\begin{chapterexercise}{宽度瓶颈}
某核心取指、译码、重命名和退休宽度均为 4，但每拍最多只有两个整数执行端口。一个长基本块全部由互不依赖的整数加法组成。忽略 miss 和分支后，稳态 IPC 上限是多少？哪个边界决定上限？
\exercisecheck{学习目标}{流水宽度由最窄的持续服务边界限制}
\end{chapterexercise}

\section*{参考解答与要点}

\begin{chapteranswer}{乱序重命名}
第一条读取 \code{p2}，为 x1 分配 \code{p20}，替代 \code{p1}；第二条读取第一条产生的 \code{p20}，为 x3 分配 \code{p21}；第三条读取 x4 的当前映射，为 x1 分配 \code{p22}，替代 \code{p20}。第二条保留对 \code{p20} 的 RAW 依赖，第三条不会覆盖它。
\answercheck{必须掌握的要点}{组内后继读取新映射；旧物理寄存器到退休后才能回收}
\end{chapteranswer}

\begin{chapteranswer}{唤醒与端口}
select 只从两条就绪操作中选择一条占用乘法端口；另一条继续留在发射队列，源就绪位保持有效，下一拍重新参与选择。清除就绪位会使它等待一个不会再次广播的结果标签。
\answercheck{必须掌握的要点}{wakeup 记录数据已到达，select 决定本拍资源归属}
\end{chapteranswer}

\begin{chapteranswer}{存储转发}
load 的四个字节完全被较早 store 覆盖，应从 store queue 转发对应高四字节。若 load 先读了 Cache，后来才发现同址更早 store，LSQ 标记内存顺序违规，使该 load 及其依赖操作失效并从安全边界重放。
\answercheck{必须掌握的要点}{转发与违规检测都按地址范围和字节掩码判断}
\end{chapteranswer}

\begin{chapteranswer}{精确异常}
提交端先让无异常的 I1、I2 按序退休；I3 到达 ROB 队首时停止提交，保存异常原因和故障地址，恢复 I3 之前的架构映射，清除 I3、I4、I5 及其队列项，然后把 PC 重定向到异常入口。I4、I5 的物理结果从未成为架构状态。
\answercheck{必须掌握的要点}{异常前已提交，异常指令及更年轻指令均不提交}
\end{chapteranswer}

\begin{chapteranswer}{误预测恢复}
至少恢复 RAT、空闲表位置、ROB 尾、发射队列、LSQ 和取指队列，并清除错误路径操作。若只恢复 RAT 而不恢复空闲表，错误路径分配的物理寄存器不会归还，运行一段时间后会出现虚假的“无空闲寄存器”停顿。
\answercheck{必须掌握的要点}{状态值和资源所有权必须回到同一检查点}
\end{chapteranswer}

\begin{chapteranswer}{宽度瓶颈}
稳态每拍最多执行两个加法，IPC 上限为 2。前端与退休虽能处理 4 条，但持续服务能力由两个整数端口限定；队列只能暂时吸收波动，不能改变长期吞吐上限。
\answercheck{必须掌握的要点}{长期吞吐取决于最窄且持续被使用的资源}
\end{chapteranswer}
